% Section 4 Relations Among Views
% Module mapping: Slice 2; deployment mapping: Slice 5 (issue #21).
% The section heading is in main.tex.
This draft contains three views: the module decomposition view --- the primary view,
Section~\ref{sec:module-view} --- the Components-and-Connectors view --- the key
runtime view, Section~\ref{sec:cc-view} --- and the Allocation view,
Section~\ref{sec:allocation-view}. The tables below relate implementation units,
runtime components and deployment nodes. Every internal runtime component
of the C\&C view is realised by exactly one module of the module view, as
Table~\ref{tab:view-mapping} records. The module view adds one module, Common, that is
implementation-only --- the shared domain types and the \IEP library --- and so has no
runtime component of its own. The \StandbyMaster is a second instance of the module that
realises the \CentralMasterInventory. The external \AnalyticsClient is outside the system:
it reads through the \InventoryQuery and realises no module of ours, so it is not mapped.

\begin{table}[htbp]
\centering
\footnotesize
\setlength{\tabcolsep}{4pt}
\renewcommand{\arraystretch}{0.95}
\begin{tabularx}{\linewidth}{|>{\raggedright\arraybackslash}p{0.30\linewidth}|>{\raggedright\arraybackslash}X|}
\hline
\textbf{Module (module view)} & \textbf{Runtime components it realises (C\&C view)}\\\hline
Store Subsystem & \POSAdapter, \StoreAgent, \StoreLocalInventory, \Outbox\\\hline
\CFE & \CFE, CFE event store\\\hline
Distribution Center Subsystem & \DCAgent, \CentralMasterInventory, \StandbyMaster, \ShippingOrderFulfillment, \StoreTrackingClient, \ReportingClient, Master inventory database\\\hline
Common & None --- shared code used by the other three modules\\\hline
\end{tabularx}
\caption{Module-to-runtime-component mapping between the module decomposition view and the Components-and-Connectors view.}
\label{tab:view-mapping}
\end{table}

Table~\ref{tab:component-node-mapping} extends that mapping with baseline deployment
locations from Figure~\ref{fig:deployment}. Each store repeats the same placements;
each listed component instance has one host. Common is shared code incorporated
into the deployed services, not an additional runtime host.

\begin{table}[htbp]
\centering
\small
\begin{tabularx}{\linewidth}{|>{\raggedright\arraybackslash}X|>{\raggedright\arraybackslash}p{0.35\linewidth}|}
\hline
\textbf{Runtime components / state} & \textbf{Execution node / storage}\\\hline
\POSAdapter, \StoreAgent & Store node, one pair per store ($\times N$)\\\hline
\StoreLocalInventory, \Outbox & Store node's local durable storage ($\times N$)\\\hline
\CFE, \DCAgent, \CentralMasterInventory, \ShippingOrderFulfillment, \StoreTrackingClient, \ReportingClient & Primary \DC node\\\hline
\StandbyMaster and its database copy & Second \DC host (Section~\ref{sec:allocation-view})\\\hline
CFE event store, master inventory database & Primary \DC node's local durable storage\\\hline
External POS system & External POS node, one per store\\\hline
\AnalyticsClient & External analytics node (future capability)\\\hline
\end{tabularx}
\caption{Component-to-node mapping from the C\&C view to the Allocation deployment view.}
\label{tab:component-node-mapping}
\end{table}

The store's recovery log/snapshot is supporting state on store-local storage;
source/subscription positions and processing checkpoints belong to the corresponding
host's durable storage. These artifacts support the recovery behavior in
Figure~\ref{fig:cc-seq-reconcile}; they add no separate application component.
The WAN connects the store and \DC hosts inside the system boundary, while POS
and analytics query links cross that boundary. Separating other \DC services onto
additional hosts requires updating this mapping and its capacity/failure assumptions.
